% ============================================================================
% PROTOCOLO EXPERIMENTAL - ENFOQUE 1: CALIDAD DE RF HUMANOS VS LLM
% Sistema AgriMoreira - Sistema de Gestion Agricola con IA
% Proyecto Fin de Curso - Ingenieria de Requerimientos [20303]
% Universidad Tecnica Estatal de Quevedo (UTEQ)
% Integrante responsable: Escudero Plaza Maria del Rosario
% Base de contenido: Aporte_Jeanpierre_v2 (titulo y resumen del manuscrito)
% Version 1.0
% ============================================================================
\documentclass[12pt,a4paper]{article}

\usepackage[utf8]{inputenc}
\usepackage[T1]{fontenc}
\usepackage[spanish,es-tabla]{babel}
\usepackage[margin=2.5cm]{geometry}
\usepackage{amsmath,amssymb}
\usepackage{booktabs}
\usepackage{graphicx}
\usepackage{hyperref}
\usepackage{longtable}
\usepackage{array}
\usepackage{multirow}
\usepackage{xcolor}

\hypersetup{
    colorlinks=true,
    linkcolor=blue,
    urlcolor=blue,
    citecolor=blue
}

\title{\textbf{Protocolo Experimental}\\[0.5em]
\large Comparacion de la calidad de requisitos funcionales elicitados por
analistas humanos frente a los generados por un modelo de lenguaje (LLM) en un
sistema agricola con IA\\[0.5em]
\large (Enfoque 1)}
\author{Escudero Plaza Maria del Rosario\\[0.3em]
\small Universidad Tecnica Estatal de Quevedo (UTEQ)\\
\small Proyecto Fin de Curso - Ingenieria de Requerimientos [20303]}
\date{Version 1.0 \textbullet\ Semestre 2026-2027}

\begin{document}

\maketitle

\begin{abstract}
El presente protocolo define el componente empirico del sistema
\textbf{AgriMoreira}, un sistema de gestion agricola con inteligencia artificial
para el apoyo a la toma de decisiones en cultivos de cacao y platanero verde.
El estudio sigue el \textbf{Enfoque 1} de la guia: comparar la calidad de los
requisitos funcionales (RF) elicitados por el equipo de analistas humanos con
los generados por un modelo de lenguaje de gran escala (LLM) a partir del mismo
material fuente (transcripcion anonimizada de una entrevista de campo). Se
aplica un diseno cuasi-experimental apareado: ambos conjuntos de RF (humano y
LLM), de un minimo de 25 requisitos cada uno, se evaluan a ciegas por
evaluadores expertos independientes en cinco dimensiones de calidad
(completitud, ausencia de ambiguedad, verificabilidad, correccion respecto de
la fuente y consistencia interna). El acuerdo inter-evaluador se mide con los
coeficientes kappa de Cohen y de Fleiss, y las puntuaciones medias se comparan
con prueba t apareada o Wilcoxon segun normalidad, con tamano del efecto y
calculo de potencia estadistica previa. El protocolo se registra de forma
previa en el Open Science Framework (OSF) \cite{nosek2018}.
\end{abstract}

\tableofcontents
\newpage

% ============================================================================
\section{Informacion general del proyecto}
% ============================================================================

\begin{center}
\begin{tabular}{>{\raggedright\arraybackslash}p{5cm} p{9cm}}
\toprule
\textbf{Elemento} & \textbf{Descripcion} \\
\midrule
Nombre del sistema & AgriMoreira: Sistema de Gestion Agricola con IA para
cultivos de verde y cacao \\
Dominio & Agroindustria (produccion de cacao y platanero verde, Manabi,
Ecuador) \\
Organizacion cliente & Agricola Moreira (seudonimo declarado en el protocolo de
anonimizacion) \\
Enfoque empirico & Enfoque 1: Calidad de RF humanos vs.\ LLM \\
Asignatura & Ingenieria de Requerimientos [20303] \\
Integrante responsable & Escudero Plaza Maria del Rosario \\
Docente supervisor & Mgs. Gleiston Ciceron Guerrero Ulloa \\
\bottomrule
\end{tabular}
\end{center}

% ============================================================================
\section{Preguntas de investigacion (formato PICOC)}
% ============================================================================

Las preguntas de investigacion se formulan siguiendo el marco PICOC
(poblacion, intervencion, comparacion, resultado, contexto)
\cite{kitchenham2007}.

\subsection{Preguntas de investigacion}

\begin{enumerate}
    \item \textbf{RQ1.} \emph{?`En que dimensiones de calidad los requisitos
    funcionales (RF) elicitados por el equipo humano difieren de los generados
    por un LLM a partir del mismo material fuente?}

    \item \textbf{RQ2.} \emph{?`Cual es el nivel de acuerdo inter-evaluador
    (kappa de Cohen y de Fleiss) entre los evaluadores expertos al puntuar
    ambos conjuntos de RF?}

    \item \textbf{RQ3.} \emph{?`Cual es la magnitud del efecto (Cohen d o
    Cliff delta) de la diferencia de calidad entre RF humanos y LLM en cada
    dimension?}
\end{enumerate}

\subsection{Aplicacion del marco PICOC}

\begin{center}
\begin{tabular}{>{\raggedright\arraybackslash}p{3cm} p{12cm}}
\toprule
\textbf{Componente} & \textbf{Descripcion} \\
\midrule
Poblacion (P) & Conjunto de requisitos funcionales del sistema AgriMoreira
(>= 25 RF humanos del ERS y >= 25 RF generados por LLM a partir del mismo
material fuente). \\
Intervencion (I) & Evaluacion a ciegas de los RF con una rubrica de cinco
dimensiones (completitud, ausencia de ambiguedad, verificabilidad, correccion
respecto de la fuente y consistencia interna), puntuando 1--5. \\
Comparacion (C) & Conjunto de RF producidos por LLM (Conjunto A) frente al
conjunto de RF elicitados por el equipo humano (Conjunto B). \\
Resultado (O) & Puntuacion media de calidad por dimension, acuerdo
inter-evaluador (kappa de Cohen y de Fleiss), tamano del efecto (d de Cohen o
delta de Cliff). \\
Contexto (C) & Desarrollo de un sistema real de gestion agricola con IA en un
dominio agroindustrial ecuatoriano, periodo 2026-2027. \\
\bottomrule
\end{tabular}
\end{center}

% ============================================================================
\section{Hipotesis}
% ============================================================================

Al tratarse de una comparacion cuantitativa, se formulan hipotesis nulas y
alternativas para cada dimension de calidad \(j \in \{1, \ldots, 5\}\):

\begin{center}
\begin{tabular}{p{13cm}}
\hline
\(H_{0j}\): No existe diferencia estadisticamente significativa entre la
puntuacion media de calidad de los RF humanos (Conjunto B) y la de los RF
generados por LLM (Conjunto A) en la dimension \(j\). \\
\hline
\(H_{1j}\): Existe una diferencia estadisticamente significativa entre ambas
puntuaciones medias en la dimension \(j\). \\
\hline
\end{tabular}
\end{center}

Nivel de significancia fijado en \(\alpha = 0{,}05\).

% ============================================================================
\section{Variables del estudio}
% ============================================================================

\subsection{Variables independientes}
\begin{itemize}
    \item Origen del conjunto de requisitos: Conjunto A (LLM) vs.\ Conjunto B
    (equipo humano).
\end{itemize}

\subsection{Variables dependientes}
\begin{itemize}
    \item Puntuacion de calidad de cada RF en cada dimension (escala 1--5):
    completitud, ausencia de ambiguedad, verificabilidad, correccion respecto
    de la fuente y consistencia interna.
    \item Acuerdo inter-evaluador: kappa de Cohen \cite{cohen1960} entre pares
    de evaluadores y kappa de Fleiss \cite{fleiss1971} para el conjunto
    completo.
\end{itemize}

\subsection{Variables de control}
\begin{itemize}
    \item Material fuente unico: la misma transcripcion anonimizada de
    entrevista para ambos conjuntos.
    \item Rubrica de evaluacion identica para ambos conjuntos y todos los
    evaluadores.
    \item Evaluadores independientes y ciegos (no saben cual conjunto es cual).
    \item Parametros del LLM fijos: modelo, temperatura, top-p y semilla
    registrados en \texttt{06\_Experimento/prompts\_llm/}.
\end{itemize}

% ============================================================================
\section{Diseño del estudio}
% ============================================================================

Cuasi-experimento apareado \cite{wohlin2012}:

\begin{enumerate}
    \item Seleccion del material fuente unico: la transcripcion completa de la
    entrevista al administrador (Entr-01, evidencia EV-01 en la matriz de
    trazabilidad, archivo \texttt{02\_Evidencias/Transcripciones/}), anonimizada.
    \item Ejecucion del LLM con la consigna exacta: \emph{``A partir del
    siguiente material fuente, redacta requisitos funcionales del sistema
    descrito, con los ocho atributos de la plantilla del silabo''}. Se registra
    la consigna completa, el modelo, la temperatura, el top-p, la semilla y la
    fecha y hora de la consulta.
    \item Conjunto A (RF generados por el LLM) y Conjunto B (RF elicitados por
    el equipo humano, tomados del ERS). Cada conjunto debe tener un minimo de
    25 requisitos individuales para permitir el analisis estadistico apareado
    \cite{montgomery2022}.
    \item Evaluacion a ciegas: tres o mas evaluadores expertos independientes
    (cinco es el numero recomendado) puntuan cada RF de 1 a 5 en las cinco
    dimensiones, sin saber cual conjunto es humano y cual es LLM.
\end{enumerate}

% ============================================================================
\section{Participantes (evaluadores expertos) y reclutamiento}
% ============================================================================

\begin{itemize}
    \item \textbf{Evaluadores expertos:} un minimo de 3 y un maximo de 5
    personas expertas independientes (por ejemplo, tres estudiantes de otro
    paralelo o tres docentes del area de Ingenieria de Requerimientos).
    \item \textbf{Criterios de inclusion:} conocimientos de ingenieria de
    requisitos y/o de especificacion de software; no haber participado en la
    elicitacion de los RF del equipo.
    \item \textbf{Criterios de exclusion:} miembros del equipo de desarrollo de
    AgriMoreira y personas que hayan redactado alguno de los conjuntos de RF.
    \item \textbf{Proceso de consentimiento:} consentimiento informado firmado,
    conforme a la Ley Organica de Proteccion de Datos Personales del Ecuador
    \cite{lopdp2021}; autorizacion explicita para el uso de datos anonimizados
    en publicaciones cientificas revisadas por pares y para el deposito del
    conjunto de datos anonimizado en Zenodo con licencia CC BY 4.0.
\end{itemize}

% ============================================================================
\section{Instrumentos}
% ============================================================================

Los instrumentos completos se encuentran en la carpeta
\texttt{06\_Experimento/instrumentos/} y los prompts en
\texttt{06\_Experimento/prompts\_llm/}:

\begin{enumerate}
    \item \textbf{Material fuente unico}: transcripcion anonimizada de la
    entrevista al administrador (Entr-01, EV-01), sin nombres propios.
    \item \textbf{Consigna exacta del LLM} con modelo, temperatura, top-p,
    semilla, fecha y hora de la consulta (archivo Markdown por experimento).
    \item \textbf{Rubrica de evaluacion experta}: cada RF se puntua de 1 a 5 en
    cinco dimensiones (completitud, ausencia de ambiguedad, verificabilidad,
    correccion respecto de la fuente y consistencia interna).
    \item \textbf{Hoja de evaluacion a ciegas}: plantilla para que los
    evaluadores registren sus puntuaciones por RF y por dimension.
    \item \textbf{Plantilla de consentimiento actualizado}: disponible en la
    carpeta etica del repositorio (\texttt{08\_Etica/}).
\end{enumerate}

% ============================================================================
\section{Plan de analisis estadistico}
% ============================================================================

\subsection{Analisis descriptivo}
Se calculan la media, la desviacion estandar y las distribuciones de
frecuencias de las puntuaciones por conjunto (A y B), por dimension y por
evaluador.

\subsection{Pruebas de hipotesis}
\begin{enumerate}
    \item Prueba de normalidad de \textbf{Shapiro-Wilk} sobre las diferencias
    apareadas.
    \item Si los datos son normales: prueba \textbf{t de Student para muestras
    apareadas}.
    \item Si no son normales: prueba \textbf{de Wilcoxon} para muestras
    apareadas.
    \item Tamano del efecto: coeficiente \textbf{d de Cohen} (datos normales)
    o \textbf{delta de Cliff} (datos no normales).
    \item Correccion por comparaciones multiples (por ejemplo, Bonferroni)
    cuando se compare mas de una dimension simultaneamente.
\end{enumerate}

\subsection{Acuerdo inter-evaluador}
\begin{itemize}
    \item kappa de Cohen \cite{cohen1960} para la concordancia entre cada par
    de evaluadores.
    \item kappa de Fleiss \cite{fleiss1971} para la concordancia del conjunto
    completo de evaluadores.
\end{itemize}

\subsection{Calculo de potencia previo}
Potencia estadistica fijada en \(1-\beta = 0{,}80\), \(\alpha = 0{,}05\), tamano
de efecto medio de Cohen \(d = 0{,}5\), conforme a la lista de verificacion de
encuestas de Molleri, Petersen y Mendes \cite{molleri2020}. El calculo se
realiza antes de la ejecucion del experimento.

\subsection{Reproducibilidad}
Todos los calculos se ejecutan con los scripts de Python versionados en
\texttt{06\_Experimento/scripts\_analisis/} a partir de los datos crudos en
\texttt{06\_Experimento/resultados/}. Ninguna tabla ni figura del manuscrito se
produce manualmente.

% ============================================================================
\section{Consideraciones eticas}
% ============================================================================

\begin{itemize}
    \item Los evaluadores y participantes firman consentimiento informado antes
    de cualquier recoleccion de datos.
    \item No se publican datos identificables: la firma, la cedula, el rostro y
    la voz permanecen en zona restringida.
    \item El proyecto es de riesgo minimo operativo (Categoria C del paquete
    etico).
    \item La adenda de la segunda ronda se encuentra en
    \texttt{08\_Etica/Adenda\_Segunda\_Ronda.pdf} y tiene fecha anterior al
    primer consentimiento de la segunda ronda.
\end{itemize}

% ============================================================================
\section{Registro previo (preregistro)}
% ============================================================================

El protocolo se registra en el Open Science Framework (\url{https://osf.io})
con fecha anterior a la ejecucion del experimento, siguiendo la revolucion del
preregistro documentada por \cite{nosek2018}. El comprobante de registro con
URL persistente se almacena como \texttt{06\_Experimento/osf\_registration.pdf}.

% ============================================================================
\section{Cronograma}
% ============================================================================

\begin{center}
\begin{tabular}{l l}
\toprule
\textbf{Actividad} & \textbf{Semana} \\
\midrule
Registro del protocolo en OSF & Semana 11 \\
Ejecucion del LLM (Conjunto A) y conformacion de los conjuntos & Semana 12 \\
Evaluacion a ciegas por los evaluadores expertos & Semana 13 \\
Analisis estadistico y figuras & Semana 14 \\
Redaccion de resultados del manuscrito & Semana 15--16 \\
\bottomrule
\end{tabular}
\end{center}

% ============================================================================
\section{Amenazas a la validez}
% ============================================================================

\begin{description}
    \item[Validez interna:] se controla que ambos conjuntos se evaluen con la
    misma rubrica y los mismos evaluadores a ciegas. Limitacion: la seleccion
    del material fuente unico puede influir en la calidad de los RF de ambos
    conjuntos.
    \item[Validez externa:] el estudio se limita a un sistema agroindustrial
    ecuatoriano y a un solo material fuente; los resultados no son
    generalizables sin estudios de replicacion.
    \item[Validez de constructo:] la calidad de requisitos se operacionaliza en
    cinco dimensiones validadas en la literatura
    \cite{montgomery2022}; se triangula con el acuerdo inter-evaluador.
    \item[Validez de conclusion:] se reportan el estadistico, los grados de
    libertad, el valor \(p\) y el tamano del efecto en cada afirmacion
    cuantitativa.
\end{description}

\bibliographystyle{unsrt}
\bibliography{referencias}

\end{document}
